\chapter{Numerical Recipes}
\label{app:numerical-recipes}

This appendix provides implementation notes for the key
computational tasks discussed in the book. Each recipe records a
reproducible algorithm for producing the corresponding figure or numerical
check. The steps are written in a language-independent form so that they
can be reimplemented with any reliable numerical package.


%% ================================================================
\section{Lorentz dispersion model}
\label{sec:recipe-lorentz}

\paragraph{Aim.}
Compute the complex permittivity $\epsr(\omega)$ for a single-oscillator
Lorentz model with passive damping or linear gain
(Figure~\ref{fig:lorentz-dispersion}).

\paragraph{Equations.}
\begin{equation}
 \epsr(\omega) = \varepsilon_\infty
 + \frac{F}{\omega_0^2 - \omega^2 - i\gamma\omega},
 \label{eq:recipe-lorentz}
\end{equation}
with the $e^{-i\omega t}$ convention. Positive $\gamma$ gives
absorption ($\Im\epsr > 0$ near resonance); negative $\gamma$ gives
gain.

\paragraph{Procedure.}
\begin{enumerate}
 \item Choose parameters: $\varepsilon_\infty$, $F$, $\omega_0$,
 $\gamma_{\mathrm{abs}} > 0$, $\gamma_{\mathrm{gain}} < 0$.
 \item Sample $\omega$ on a uniform grid in
 $[0.15\,\omega_0,\; 2.25\,\omega_0]$.
 \item Evaluate~\eqref{eq:recipe-lorentz} for both $\gamma$ values.
 \item Plot $\Re\epsr(\omega)$ and $\Im\epsr(\omega)$ on separate
 panels.
\end{enumerate}

\paragraph{Result.}
This calculation gives the Lorentz-dispersion figure
\texttt{fig\_f5\_lorentz\_dispersion.pdf}.


%% ================================================================
\section{Two-mode exceptional point}
\label{sec:recipe-ep}

\paragraph{Aim.}
Compute eigenvalue surfaces and eigenvector coalescence for the
canonical $2\times 2$ EP Hamiltonian
(Figures~\ref{fig:ep-eigenvalue-surfaces}
and~\ref{fig:ep-eigenvector-coalescence}).

\paragraph{Equations.}
\begin{equation}
 H = \begin{pmatrix} i\gamma & \kappa \\ \kappa & -i\gamma
 \end{pmatrix}, \qquad
 \lambda_\pm = \pm\sqrt{\kappa^2 - \gamma^2}.
 \label{eq:recipe-ep}
\end{equation}

\paragraph{Procedure.}
\begin{enumerate}
 \item Create a 2D grid in $(\kappa, \gamma)$.
 \item Evaluate $\lambda_\pm = \pm\sqrt{\kappa^2 - \gamma^2}$ on the
 grid (complex square root).
 \item Plot $\Re\lambda_\pm$ and $\Im\lambda_\pm$ as 3D surfaces.
 \item For the eigenvector figure: fix $\kappa = 1$, sweep $\gamma$,
 diagonalize $H$ numerically, compute the overlap
 $|\langle v_+|v_-\rangle|$ and the condition number of the
 eigenvector matrix.
 \item For the encircling figure: parameterize a loop
 $\delta = r\,e^{i\theta}$ around the branch point, plot
 $\pm\sqrt{\delta}$ in the complex $\lambda$ plane.
\end{enumerate}

\paragraph{Result.}
This calculation gives the EP eigenvalue and eigenvector figures
\texttt{fig\_f8\_ep\_eigenvalues.pdf} and
\texttt{fig\_f8\_ep\_vectors.pdf}.


%% ================================================================
\section{Point-gap winding number}
\label{sec:recipe-winding}

\paragraph{Aim.}
Compute the spectral winding number of a complex band loop around a
reference energy (see the complex-band-loop discussion associated with
Figure~\ref{fig:complex-band-loop}).

\paragraph{Equations.}
One-band non-Hermitian dispersion:
\begin{equation}
 \omega(k) = \omega_0 + t_R\,e^{ik} + t_L\,e^{-ik} + i\gamma_0,
 \label{eq:recipe-winding}
\end{equation}
with $t_R \neq t_L$ producing a loop that can encircle the origin.

\paragraph{Algorithm.}
\begin{enumerate}
 \item Sample $k$ uniformly on $[-\pi, \pi]$ with $N \gg 1$ points.
 \item Evaluate $\omega(k)$ and form $z(k) = \omega(k) - E_{\mathrm{ref}}$.
 \item Compute the unwrapped phase:
 $\phi(k) = \mathrm{unwrap}[\arg z(k)]$.
 \item The winding number is $W = [\phi(\pi) - \phi(-\pi)]/(2\pi)$.
 \item Compare inside and outside reference points: an inside point
 gives $|W| \ge 1$; an outside point gives $W = 0$.
\end{enumerate}

\paragraph{Result.}
The calculation gives a point-gap winding diagnostic and the
complex-band-loop figure \texttt{fig\_f13\_complex\_band\_loop.pdf}.
If the standalone point-gap-winding schematic is omitted from the main
text, the same calculation can be reported directly in the discussion of
the remaining loop figure.


%% ================================================================
\section{Non-Hermitian skin effect (Hatano--Nelson model)}
\label{sec:recipe-skin}

\paragraph{Aim.}
Demonstrate the PBC/OBC spectral discrepancy, eigenmode localization,
and the generalized Brillouin zone for a minimal skin-effect model
(Figures~\ref{fig:skin-pbc-obc-spectra},
\ref{fig:skin-mode-localization}, and~\ref{fig:gbz-circle}).

\paragraph{Equations.}
Hatano--Nelson chain: $H_{n,n+1} = t_R$, $H_{n+1,n} = t_L$, with
$t_R > t_L$.

\paragraph{Algorithm.}
\begin{enumerate}
 \item \textbf{PBC spectrum:} $E_{\mathrm{PBC}}(k)
 = t_R\,e^{ik} + t_L\,e^{-ik}$; plot as a complex loop.
 \item \textbf{OBC spectrum:} construct the $N\times N$ tridiagonal
 matrix, diagonalize, plot eigenvalues in the complex plane.
 \item \textbf{Mode localization:} normalize $|\psi_j(n)|^2$ for each
 eigenstate; compute center of mass
 $\bar{n}_j = \sum_n n\,|\psi_j(n)|^2$.
 \item \textbf{GBZ:} the GBZ radius is
 $r = \sqrt{t_L/t_R}$; plot the circle $|\beta| = r$ and compare
 $E_{\mathrm{GBZ}}(\theta) = t_R\,r\,e^{i\theta}
 + t_L\,e^{-i\theta}/r$ with the OBC spectrum.
\end{enumerate}

\paragraph{Result.}
The calculation gives the PBC/OBC spectra, skin-mode localization, and
GBZ figures: \texttt{fig\_f14\_pbc\_obc\_spectra.pdf},
\texttt{fig\_f14\_skin\_modes.pdf}, and \texttt{fig\_f14\_gbz.pdf}.


%% ================================================================
\section{BIC topological charge and Q scaling}
\label{sec:recipe-bic}

\paragraph{Aim.}
Visualize the far-field polarization winding around a BIC and the
quality-factor scaling under charge merging
(Figures~\ref{fig:bic-polarization-charge}
and~\ref{fig:bic-q-scaling}).

\paragraph{Algorithm.}
\begin{enumerate}
 \item \textbf{Polarization field:} on a 2D grid in $(k_x, k_y)$,
 set the polarization angle $\theta = q\,\mathrm{atan2}(k_y, k_x)$
 with $q = +1$ or $q = -1$. Plot as a quiver field; the BIC is the
 singularity at the origin.
 \item \textbf{Q scaling:} for an isolated charge-one BIC,
 $\gamma \propto k^2$ so $Q \propto k^{-2}$. For a merged super-BIC,
 $\gamma \propto k^6$ so $Q \propto k^{-6}$. Plot both on a
 log--log scale.
 \item \textbf{BZF scaling:} $Q \sim 1/(\alpha^2 k^2)$ for
 perturbation strength $\alpha$; plot several $\alpha$ values.
\end{enumerate}

\paragraph{Result.}
The calculation gives the BIC polarization-charge and quality-factor scaling
figures, \texttt{fig\_f16\_bic\_polarization\_charge.pdf} and
\texttt{fig\_f16\_q\_scaling.pdf}.


%% ================================================================
\section{PML validation and mode-ratio filtering}
\label{sec:recipe-pml}

\paragraph{Aim.}
Illustrate the PML validation protocol: physical QNM poles remain stable
under PML parameter variation, while PML continuum modes move
(Figures~\ref{fig:pml-convergence} and~\ref{fig:pml-mode-ratio}).

\paragraph{Algorithm.}
\begin{enumerate}
 \item Generate synthetic eigenfrequencies: one physical QNM near
 a target pole, plus several PML modes whose positions depend
 strongly on the PML absorption strength $\sigma_{\max}$.
 \item Sweep $\sigma_{\max}$; for each value, compute the mode ratio
 $\mathrm{MR} = \int_{\Omega_{\mathrm{phys}}}|\EE|^2\,\dd V
 \big/ \int_{\Omega_{\mathrm{total}}}|\EE|^2\,\dd V$.
 \item Physical QNMs cluster near $\mathrm{MR} \approx 1$; PML modes
 near $\mathrm{MR} \ll 1$.
 \item Plot the eigenvalue trajectories and the mode-ratio
 separation.
\end{enumerate}

\paragraph{Result.}
The calculation gives the PML-convergence and PML-mode-ratio figures,
\texttt{fig\_f17\_pml\_convergence.pdf} and
\texttt{fig\_f17\_pml\_mode\_ratio.pdf}.


%% ================================================================
\section{Scattering pole-zero map}
\label{sec:recipe-pole-zero}

\paragraph{Aim.}
Visualize the S-matrix pole and zero trajectories under gain/loss
tuning, the CPA-laser threshold, and the complex scattering response
landscape discussed in Chapter~\ref{ch:poles-zeros}.

\paragraph{Equations.}
Scalar scattering response:
\begin{equation}
 s(\omega; g) = \prod_j \frac{\omega - z_j(g)}{\omega - p_j(g)},
 \label{eq:recipe-pole-zero}
\end{equation}
where poles $p_j$ rise from $\Im\omega < 0$ and zeros $z_j$ descend
from $\Im\omega > 0$ as the gain parameter $g$ increases.

\paragraph{Algorithm.}
\begin{enumerate}
 \item Define two pole--zero pairs near a central frequency, plus one
 detuned background pair.
 \item Sweep $g$ from $0$ to above threshold ($g = 1$).
 \item Plot pole and zero trajectories; mark the real-axis crossing
 (lasing/CPA/CPA-laser).
 \item Evaluate $|s(\omega)|^2$ on the real axis for several $g$
 values; plot spectra.
 \item Evaluate $\log_{10}|s(\omega)|$ on a 2D complex-$\omega$ grid
 near threshold; display as a heatmap.
\end{enumerate}

\paragraph{Result.}
The calculation gives pole--zero trajectories and scattering-response
data. If the schematic pole--zero figures are not included in the main
text, the recipe remains useful as a numerical check on the meromorphic
structure of the scattering amplitude.


%% ================================================================
\section{EP sensing response versus noise}
\label{sec:recipe-ep-sensing}

\paragraph{Aim.}
Compare the frequency splitting, linewidth, and signal-to-noise ratio
of an EP sensor with a conventional diabolic-point sensor, following the
scaling discussion in Chapter~\ref{ch:devices}.

\paragraph{Equations.}
Diabolic-point splitting: $\Delta\omega_{\mathrm{DP}} \propto \epsilon$.
EP splitting: $\Delta\omega_{\mathrm{EP}} \propto \sqrt{\epsilon}$.
Noise proxy: $\sigma_{\mathrm{EP}} \propto \sqrt{1 + c/\epsilon}$
(diverges as $\epsilon \to 0$).

\paragraph{Algorithm.}
\begin{enumerate}
 \item Sample $\epsilon$ on a logarithmic grid.
 \item Compute signal, noise, responsivity, and SNR for both sensor
 types.
 \item Plot three panels: splitting, responsivity+noise, and SNR
 comparison.
\end{enumerate}

\paragraph{Result.}
The calculation gives the EP sensing response-versus-noise diagnostic;
depending on the presentation, it may be shown as a figure or summarized
as the scaling argument in Chapter~\ref{ch:devices}.


%% ================================================================
\section{Practical convergence protocol}
\label{sec:convergence-protocol}

This section collects a systematic convergence protocol for
non-Hermitian photonic eigenvalue calculations. Each step can be
implemented as a post-processing check on the output of any FEM or
FDTD solver.

\subsection{Mesh refinement}
\label{ssec:mesh-convergence}

The first convergence check is mesh refinement. Run the eigenvalue
calculation at two or more mesh densities (e.g.\ $h$, $h/2$, $h/4$)
and monitor the change in the complex eigenfrequency
$\tilde{\omega}_n$. For a well-resolved mode, the eigenfrequency
should converge as
\begin{equation}
 |\tilde{\omega}_n(h) - \tilde{\omega}_n(h/2)|
 \sim C\,h^{2p},
 \label{eq:mesh-convergence}
\end{equation}
where $p$ is the polynomial order of the finite-element basis. If
the convergence rate is slower than expected, the mode may be
under-resolved---common for tightly confined plasmonic modes or modes
with sharp field gradients at material interfaces.

\paragraph{Recommended criterion.}
Require $|\Delta\tilde{\omega}| / |\tilde{\omega}| < 10^{-4}$ between
the two finest meshes. For quality factors exceeding $10^4$, tighter
convergence ($< 10^{-6}$) may be needed because $Q = \Re\tilde{\omega}
/(2|\Im\tilde{\omega}|)$ amplifies small errors in
$\Im\tilde{\omega}$.

\subsection{PML thickness and absorption profile}
\label{ssec:pml-convergence}

The PML is an artificial absorbing boundary that must be thick enough
to suppress reflections but thin enough to keep the computational
domain manageable. The PML convergence test varies the PML thickness
$d_{\mathrm{PML}}$ (or equivalently the PML absorption strength
$\sigma_{\max}$) while keeping the physical domain fixed.

\paragraph{PML thickness selection rule.}
For a mode with wavelength $\lambda$ and exterior decay constant
$\kappa_{\mathrm{out}}$, a practical starting point is
\begin{equation}
 d_{\mathrm{PML}} \gtrsim
 \max\left(\lambda,\frac{4}{\kappa_{\mathrm{out}}}\right),
 \label{eq:pml-thickness-rule}
\end{equation}
followed by an explicit sweep of thickness and absorption strength.
This estimate is only a starting scale: radiative QNMs, grazing waves,
and quasi-BICs often require thicker PMLs and finer PML discretization
than any one-line rule can predict.

\paragraph{PML eigenvalue filtering.}
Not all eigenvalues returned by the solver correspond to physical QNMs.
The PML itself introduces a branch of spurious ``PML modes'' whose
eigenvalues depend on the PML parameters. These are identified by
their sensitivity to PML changes: physical QNM eigenvalues remain
stable as $d_{\mathrm{PML}}$ varies, while PML modes drift
systematically. A mode-ratio filter
\begin{equation}
 \mathrm{MR}_n = \frac{\int_{\mathrm{physical}} |\EE_n|^2\;\dd V}
 {\int_{\mathrm{total}} |\EE_n|^2\;\dd V}
 \label{eq:mode-ratio-filter}
\end{equation}
provides a second diagnostic: physical QNMs have $\mathrm{MR} \approx 1$,
while PML modes have $\mathrm{MR} \ll 1$. In practice, a threshold
$\mathrm{MR} > 0.5$ is sufficient to filter out most PML artifacts.

\subsection{Dispersive material branch verification}
\label{ssec:dispersion-branches}

For dispersive materials with multiple resonances (e.g.\ Lorentz
oscillators), the analytic continuation $\epsr(\omega)$ to complex
$\omega$ must use the same causal branch as the real-frequency material
model. Eigenvalues that are artifacts of a wrong analytic continuation
are unphysical. Useful checks are:
\begin{enumerate}
 \item Evaluate $\epsr(\tilde{\omega}_n)$ at the computed complex
 eigenfrequency.
 \item Track the mode continuously from a real-frequency or weakly open
 reference calculation while keeping the material poles and square-root
 branches fixed.
 \item Check passivity on the real axis for the underlying material
 model; the sign of $\Im\epsr(\tilde{\omega}_n)$ at a complex
 eigenfrequency is not, by itself, a reliable physicality test.
 \item Discard eigenvalues that switch sheets, converge to material
 poles rather than electromagnetic resonances, or fail the residual and
 PML-variation checks.
\end{enumerate}

\subsection{Biorthogonal normalisation check}
\label{ssec:biorthogonal-check}

After extracting the QNM eigenvectors, verify the biorthogonal
normalisation condition
\begin{equation}
 \langle \tilde{\EE}_m, \tilde{\EE}_n \rangle_{\mathrm{uc}}
 = \delta_{mn},
 \label{eq:biorthogonal-check}
\end{equation}
where $\langle\cdot,\cdot\rangle_{\mathrm{uc}}$ is the unconjugated
inner product (Chapter~\ref{ch:qnm}). Deviations from $\delta_{mn}$
indicate either insufficient mesh resolution, incorrect PML
regularisation, or missing modes. A practical tolerance is
$|\langle \tilde{\EE}_m, \tilde{\EE}_n \rangle_{\mathrm{uc}} -
\delta_{mn}| < 10^{-3}$.

\subsection{Reproducibility checklist}
\label{ssec:reproducibility}

For publication-quality results, record the following for every
eigenvalue calculation:
\begin{center}
\begin{tabular}{ll}
 \hline
 \textbf{Parameter} & \textbf{Purpose} \\
 \hline
 Mesh element size $h$ & Resolution verification \\
 FE polynomial order $p$ & Convergence rate check \\
 PML thickness $d_{\mathrm{PML}}$ & Artificial-boundary control \\
 PML absorption $\sigma_{\max}$ & Reflection suppression \\
 Dispersion model & Material branch verification \\
 Contour parameters ($\mathcal{C}$, $N_q$) & Eigensolver control \\
 Mode-ratio threshold & PML filtering \\
 Biorthogonal norm tolerance & Normalisation quality \\
 \hline
\end{tabular}
\end{center}
A compact reporting table of this kind is recommended for numerical QNM
results that support physical claims. Its consistent use would
significantly improve the reproducibility of non-Hermitian photonic
calculations.



\subsection{Minimum reporting standard for Maxwell eigenvalue computations}
\label{ssec:min-reporting-standard}

A non-Hermitian Maxwell calculation is not fully reproducible unless the
reported eigenfrequency is accompanied by an estimate of its numerical
uncertainty and by a check that the computed left and right modes solve
the same physical problem.  The following minimal standard is recommended
for any result that is used to support a physical claim.

First, report the residual norm of the right eigenproblem.  Consider a
discretised generalized Maxwell problem evaluated at the nonlinear
eigenfrequency $\tilde{\omega}_n$:
\begin{equation}
  \mathbf{A}(\tilde{\omega}_n)\mathbf{x}_n
  = \tilde{\omega}_n^2\,\mathbf{B}(\tilde{\omega}_n)\mathbf{x}_n,
  \label{eq:reporting-gep}
\end{equation}
the dimensionless residual may be written as
\begin{equation}
  \rho_n^{\mathrm{R}}
  = \frac{\|\mathbf{A}(\tilde{\omega}_n)\mathbf{x}_n
  - \tilde{\omega}_n^2\mathbf{B}(\tilde{\omega}_n)\mathbf{x}_n\|_2}
  {\|\mathbf{A}(\tilde{\omega}_n)\mathbf{x}_n\|_2
  + |\tilde{\omega}_n|^2\|\mathbf{B}(\tilde{\omega}_n)\mathbf{x}_n\|_2}.
  \label{eq:right-residual}
\end{equation}
For a well-converged mode, $\rho_n^{\mathrm{R}}$ should be much smaller
than the quoted relative uncertainty in $\tilde{\omega}_n$.  Reporting
only the eigenfrequency without the residual is insufficient, because
spurious PML modes and poorly resolved physical modes can have plausible
complex frequencies.

Second, when using biorthogonal perturbation theory, report the residual
of the left eigenproblem as well.  If the adjoint mode is represented by
$\mathbf{y}_n$, the corresponding residual is
\begin{equation}
  \rho_n^{\mathrm{L}}
  = \frac{\|\mathbf{A}^{\mathrm{T}}(\tilde{\omega}_n)\mathbf{y}_n
  - \tilde{\omega}_n^2\mathbf{B}^{\mathrm{T}}(\tilde{\omega}_n)\mathbf{y}_n\|_2}
  {\|\mathbf{A}^{\mathrm{T}}(\tilde{\omega}_n)\mathbf{y}_n\|_2
  + |\tilde{\omega}_n|^2\|\mathbf{B}^{\mathrm{T}}(\tilde{\omega}_n)\mathbf{y}_n\|_2}.
  \label{eq:left-residual}
\end{equation}
The transpose, rather than the Hermitian conjugate, appears because the
reciprocal Maxwell problem uses the unconjugated bilinear pairing.  If a
paper reports Petermann factors, adjoint sensitivities, or mode-volume
corrections but gives no left-mode residual, the numerical evidence is
incomplete.

Third, quote an uncertainty band obtained from at least two independent
refinement directions.  Mesh refinement tests the spatial discretisation;
PML variation tests the artificial boundary; material-branch checks test
the analytic continuation.  A conservative uncertainty estimate is
\begin{equation}
  \Delta\tilde{\omega}_n
  = \max\left\{
  |\tilde{\omega}_n(h)-\tilde{\omega}_n(h/2)|,
  |\tilde{\omega}_n(d_{\mathrm{PML}})-\tilde{\omega}_n(2d_{\mathrm{PML}})|,
  |\tilde{\omega}_n(N_q)-\tilde{\omega}_n(2N_q)|
  \right\},
  \label{eq:qnm-uncertainty-band}
\end{equation}
where $N_q$ is the number of quadrature points in a contour eigensolver
when such a method is used.  The reported real frequency, linewidth, and
quality factor should be rounded consistently with this uncertainty.
For example, quoting $Q=1.23748\times10^5$ is misleading if the PML
sweep changes $Q$ by three percent.

Finally, state which modes were discarded and why.  Non-Hermitian
Maxwell spectra contain physical resonances, PML continuum modes,
material-pole artifacts, and numerical outliers.  A transparent report
should say how each class was separated.  The combination of residual
norms, mode-ratio filtering, branch verification, and refinement sweeps
is usually sufficient.  Without these checks, an apparently impressive
exceptional point or ultra-high-$Q$ resonance may simply be a numerical
artifact of the discretised open boundary.
